\chapter{Market: TAM, SAM and SOM}\label{ch:market}

\section{Method and its limits}\label{sec:market-method}

We size the market bottom-up, in orders and in the fee that ClosetTwin would earn on them, and we call the result a \emph{fee pool} and not a market size. A fee pool is the number of orders a segment could send through the service multiplied by our blended price, so it measures what we could earn and not what retailers or shoppers spend. Each input carries one of three marks: sourced, derived by us from other inputs, or assumed by us with a stated basis (\ctlegend{}). Table~\ref{tab:sam-inputs} lists them all.

The Israeli figures need a warning. The market numbers come from Statista estimates that a research report of the Knesset Research and Information Center relays \cite{knesset_ric2025}. They are not a release of the Central Bureau of Statistics, the 2024 value is an estimate, and the fashion share is the 2022 value. Publisher estimates of the virtual-fitting-room market are context only and enter no calculation: IMARC Group reports that this global market reached \$8.5 billion in 2025 and will reach \$52.6 billion by 2034, a compound annual growth rate of 21.80\% \cite{imarc2026}, but the definition and method sit in a paid report that we could not read, and we found no second publisher that we could verify.

\section{Israeli online commerce in context}\label{sec:market-context}

Israeli online commerce is small next to the markets in which fitting vendors usually operate, but it is growing. The Knesset report, citing Statista, puts total online trade at about \$6.1 billion in 2023, or about NIS~22.6 billion, which is about 2.6\% of household private consumption; the estimate for 2024 is \$6.9 billion and the forecast for 2029 is \$11.3 billion \cite{knesset_ric2025}. The ratio of the two 2023 figures gives a rate of 3.70 shekels to the dollar, and we use this derived rate whenever a shekel equivalent is shown, for example \$0.036 per order, which is about \ils{0.13}. Fashion is the largest online category: in 2022 it was about 21.6\% of Israeli online commerce, ahead of electronics at about 20\% \cite{knesset_ric2025}.

\section{TAM: Israeli online fashion}\label{sec:tam}

For the total addressable market we take the 2022 fashion share of 21.6\% and hold it at its 2022 value for 2024, which gives \derived{\$1.49 billion} of online fashion sales in Israel, about NIS~5.5 billion. We multiply the share and not the shekel figure for 2022 because the report's own 2022 figures imply a rate of 3.37 shekels to the dollar, and a 2022 shekel value converted at the 2023 rate would mix two years. At an assumed average order value of \assumed{\$80} this is \derived{18.6 million} online fashion orders a year. At the blended price of \$0.036 per order, the fee pool of the whole Israeli online fashion market is only \derived{\$0.67 million} a year.

\section{SAM: local Israeli e-tailers only}\label{sec:sam}

The serviceable market is smaller still. Foreign sites are prominent among the sellers that Israeli shoppers use. The 2020 sales ranking in the Knesset report is led by Amazon with \$195 million, followed by Shein (\$151 million), and it also lists Next (\$96 million), ASOS (\$57 million) and Zara (\$51 million) \cite{knesset_ric2025}. The same report adds two charts by share of orders from November 2024, which it takes from Statista. Among global sites, AliExpress leads with 57\%, followed by Shein (12\%), Amazon and Temu (7\% each) and Next (5\%); among local sites, Super-Pharm leads with 22\%, followed by Shufersal (20\%), KSP (15\%) and Terminal X (12\%) \cite{knesset_ric2025}. The charts are two separate rankings and the report does not classify the sites by category, so they do not give a split of fashion orders between local and foreign sites. A foreign site that serves Israeli shoppers will not buy a plug-in from an Israeli start-up for that market. We therefore assume that \assumed{45\%} of Israeli online fashion orders go to local e-tailers, which is \derived{8.4 million} orders. The order shares above name individual local sites but give no number of orders for any Israeli fashion e-tailer, so we assume three size bands: 5 large retailers with 600,000 orders a year each, 25 mid-size retailers with 100,000 each and 250 small shops with 12,000 each. The bands add up to 8.5 million orders, within 2\% of the 8.4 million above. We address the large and mid-size retailers only, since a shop with 12,000 orders pays about \$430 a year and cannot justify an integration. That leaves 30 candidate retailers and \derived{5.5 million} orders, a fee pool of \derived{\$198,000} a year.

\begin{insightbox}[Israel is a reference market, not a revenue market]
The Israeli SAM fee pool is about \$0.2 million a year. No team can live on that. The Israeli plan is therefore a strategy for reference customers, for pilots that can prove the effect, and for early learning. Revenue must come from Europe.
\end{insightbox}

Why start in a market this small? There are three reasons, and none of them is revenue. First, the plan's central hypothesis can be tested only with a real retailer, and the five large Israeli retailers of our bands have enough orders for an individual test (Chapter~\ref{ch:gtm}). Second, a pilot is cheap: we assume a cash cost of \assumed{\$3,000} for a direct pilot (Chapter~\ref{ch:finance}). Third, Israeli retailers work under a statutory 14-day cancellation right (Chapter~\ref{ch:intro}), so returns are a real cost to them. The price of this choice is stated openly: by the model, the cumulative operating result at the end of year~3, before Europe opens, is a deficit of \derived{\$1.04 million}, spent in a market whose whole SAM is about \$0.2 million a year. A reference customer in Israel is useful in Europe only if European retailers accept it, and we have no evidence of that. The model does not include a route to Europe through Israeli retailers that sell to EU customers; its European clients are European retailers that the sales team reaches directly.

\begin{table}
\caption{Inputs of the market sizing, with the class of each}\label{tab:sam-inputs}
\begin{tabularx}{\textwidth}{lrlX}
\toprule
\ctheaderrow\cthead{Input} & \cthead{Value} & \cthead{Class} & \cthead{Source or basis}\\
\midrule
Online commerce 2024 & \$6.9\,B & sourced & Knesset report \cite{knesset_ric2025}\\
Fashion share, 2022 & 21.6\% & sourced & same report\\
Online fashion sales 2024 & \$1.49\,B & derived & share $\times$ \$6.9\,B\\
Average order value & \$80 & assumed & mid-market basket\\
Online fashion orders & 18.6\,M & derived & sales $\div$ order value\\
Share going to local e-tailers & 45\% & assumed & largest sites are foreign\\
Retailers (large / mid / small) & 5 / 25 / 250 & assumed & bands; order shares only in the report\\
SAM orders (large and mid) & 5.5\,M & derived & 5 $\times$ 600,000 + 25 $\times$ 100,000\\
Blended price per order & \$0.036 & derived & 60\% at \$0.04, 40\% at \$0.03\\
Israeli SAM fee pool & \$198,000 & derived & SAM orders $\times$ price\\
Europe share of global apparel & 25\% & assumed & non-APAC is 57.3\%\\
Global e-commerce apparel 2024 & \$726\,B & sourced & Kings Research \cite{kings2026}\\
Europe fee pool (TAM / SAM) & \$81.7\,M / \$32.7\,M & derived & orders $\times$ price; SAM at 40\%\\
\bottomrule
\end{tabularx}
\end{table}

\begin{figure}
% F1: TAM / SAM / SOM, Israel (nested) and Europe layer (separate). Not to scale.
% No caption/label here: the chapter owns them (fig:f1).
% Numbers: python/model.py D30, D34b_tam, D33, D34, SUBI_5, D35, D36, D37, SUBE_5. D = derived.
\begin{tikzpicture}[x=0.37cm,y=0.4cm]
  \node[ct title,anchor=south west] at (0,16.1) {Israel (beachhead)};
  \draw[ct node ivory] (0,0) rectangle (17.5,16);
  \node[ct label,anchor=north west,align=left] at (0.4,15.7)
    {\textbf{TAM} online fashion: 18.6\,M orders,\\fee pool \$0.67\,M a year\textsuperscript{D}};
  \draw[ct node pink,minimum width=0pt] (1,1) rectangle (16.5,10.8);
  \node[ct label,anchor=north west,align=left] at (1.4,10.5)
    {\textbf{SAM} local e-tailers, large and mid:\\5.5\,M orders, fee pool \$0.20\,M a year\textsuperscript{D}};
  \draw[ct key,minimum width=0pt] (2.3,2) rectangle (15.2,5.8);
  \node[ct label,text=ivory,align=center] at (8.75,3.9)
    {\textbf{SOM}, year 5\\Israeli subscription \$54\,k\textsuperscript{D}};

  \node[ct title,anchor=south west] at (20,16.1) {Europe layer (from year 4)};
  \draw[ct node ivory,dashed] (20,5.5) rectangle (37.5,16);
  \node[ct label,anchor=north west,align=left] at (20.4,15.7)
    {\textbf{TAM} 2.27\,B orders,\\fee pool \$81.7\,M a year\textsuperscript{D}};
  \draw[ct node pink,dashed,minimum width=0pt] (21,6.5) rectangle (36.5,11.2);
  \node[ct label,anchor=north west,align=left] at (21.4,10.9)
    {\textbf{SAM} fee pool\\\$32.7\,M a year\textsuperscript{D}};
  \draw[ct key,minimum width=0pt] (21,0.5) rectangle (36.5,4.5);
  \node[ct label,text=ivory,align=center] at (28.75,2.5)
    {\textbf{SOM}, year 5: subscription \$0.32\,M\textsuperscript{D}\\(not in the Y1--Y3 base case)};
\end{tikzpicture}

\caption{Market layers, not to scale: Israel (nested TAM, SAM, SOM) and the Europe layer; fee pools per year}\label{fig:f1}
\end{figure}

\section{Europe layer and SOM}\label{sec:europe-som}

Because the Israeli SAM cannot fund a team, Israel is the proof market, and Europe, entered in year 4, carries the revenue (Figure~\ref{fig:f1}). Kings Research values the global e-commerce apparel market at \$726.05 billion in 2024, of which Asia-Pacific holds 42.70\% \cite{kings2026}. The remaining 57.3\% is divided between North America, Europe and the rest of the world in a way that no source we hold gives, so we assume that Europe holds \assumed{25\%} of the global total. Divided by the same \$80 order, this is \derived{2.27 billion} orders and, at the blended price, a Europe fee pool of \derived{\$81.7 million} a year. Applying the same 40\% size filter that excludes marketplaces and small shops, which we assume for Europe only, leaves a Europe SAM of \derived{\$32.7 million}. Together with the Israeli SAM, the combined serviceable pool is \$32.9 million a year. Europe is not part of the base case of the first three years, and the plan does not claim revenue there before year~5.

The serviceable obtainable market (SOM) is the part of that pool the plan expects to win by year~5. Total revenue in year~5 is \derived{\$782,000}. Of this, Israeli subscription revenue is only \$54,000, European subscription revenue is \$648,000 and pilot fees make up the rest (Chapter~\ref{ch:finance}). Subscription revenue alone is 2.1\% of the combined SAM. This is a modest share, and the plan deliberately asks for a small slice of a pool that is still large enough to support a company. The weak link is not the size of the share but the assumption behind it, that half of the pilots become paying clients; Chapter~\ref{ch:gtm} says how the pilots test it.

\ctchapterend
